\textbf{(i)} implies \textbf{(ii)} because an isomorphism reflects units. Conversely, (ii) gives a map $T^{-1}A\to S^{-1}A$ by the universal property of localization. Its composites with $\varphi$ are the identity, again by uniqueness, proving (i).

If $t/1$ has inverse $a/s$, some $u\in S$ satisfies $u(ta-s)=0$, so $(ua)t=us\in S$. Thus \textbf{(ii)} implies \textbf{(iii)}. Conversely, $xt\in S$ makes $x/(xt)$ an inverse of $t/1$.

The criterion $\overline S=\{t:xt\in S\text{ for some }x\in A\}$ from \exref{3.7} proves \textbf{(iii)} $\Leftrightarrow$ \textbf{(iv)}. The description of $A\setminus\overline S$ as the union of primes disjoint from $S$ proves \textbf{(iv)} $\Leftrightarrow$ \textbf{(v)}.
